\documentclass[aspectratio=169,11pt]{beamer}

% ============================================================
% PRUEBAS DE HIPÓTESIS
% Curso: Ciencia de Datos para Biólogos
% Duración: 2 horas
% ============================================================

\usetheme{Madrid}
\usecolortheme{seahorse}

\usepackage[utf8]{inputenc}
\usepackage[T1]{fontenc}
\usepackage[spanish,es-nodecimaldot]{babel}
\usepackage{amsmath,amssymb}
\usepackage{booktabs}
\usepackage{graphicx}
\usepackage{array}
\usepackage{xcolor}
\usepackage{tikz}
\usepackage{listings}

% ------------------------------------------------------------
% Colores
% ------------------------------------------------------------
\definecolor{bioGreen}{RGB}{38,119,74}
\definecolor{bioBlue}{RGB}{42,91,140}
\definecolor{bioRed}{RGB}{170,55,55}
\definecolor{bioGray}{RGB}{90,90,90}

\setbeamercolor{structure}{fg=bioGreen}
\setbeamercolor{frametitle}{fg=bioGreen}
\setbeamercolor{title}{fg=bioGreen}

% ------------------------------------------------------------
% Listings
% ------------------------------------------------------------
\lstset{
  basicstyle=\ttfamily\small,
  keywordstyle=\color{bioBlue},
  commentstyle=\color{bioGray},
  stringstyle=\color{bioRed},
  showstringspaces=false,
  frame=single,
  breaklines=true
}

% ------------------------------------------------------------
% Información
% ------------------------------------------------------------
\title[Pruebas de hipótesis]{Pruebas de hipótesis}
\subtitle{Ciencia de Datos para Biólogos}
\author{Curso de Ciencia de Datos}
\institute{Biología / Ciencias Naturales}
\date{}

% ============================================================
\begin{document}
% ============================================================

% ------------------------------------------------------------
\begin{frame}
  \titlepage
\end{frame}

% ------------------------------------------------------------
\begin{frame}{Objetivos de aprendizaje}

Al finalizar la clase, el estudiante podrá:

\begin{itemize}
    \item Formular hipótesis estadísticas a partir de preguntas biológicas.
    \item Diferenciar hipótesis nula y alternativa.
    \item Interpretar un valor-$p$ correctamente.
    \item Explicar los errores tipo I y tipo II.
    \item Seleccionar pruebas de hipótesis apropiadas para distintos diseños.
    \item Distinguir entre significancia estadística y relevancia biológica.
    \item Reportar resultados de una prueba de hipótesis.
\end{itemize}

\end{frame}

% ------------------------------------------------------------
\begin{frame}{Ruta de la clase}

\begin{enumerate}
    \item Motivación: preguntas biológicas
    \item ¿Qué es una hipótesis estadística?
    \item Hipótesis nula y alternativa
    \item Estadístico de prueba
    \item Valor-$p$
    \item Errores tipo I y II
    \item Potencia estadística
    \item Pruebas frecuentes en biología
    \item Supuestos y alternativas no paramétricas
    \item Ejemplos con datos biológicos
    \item Buenas prácticas y errores frecuentes
    \item Ejercicio integrador
\end{enumerate}

\vspace{0.3cm}

\centering
\textbf{Idea central:} una prueba de hipótesis es una herramienta para
cuantificar
la evidencia de los datos frente a una hipótesis.

\end{frame}

% ------------------------------------------------------------
\begin{frame}{¿Por qué necesitamos pruebas de hipótesis?}

Supongamos que estudiamos una especie de planta.

\begin{block}{Pregunta biológica}
¿La aplicación de un fertilizante aumenta el crecimiento de las plantas?
\end{block}

Observamos:

\[
\bar{x}_{fertilizante}=18.4\text{ cm}
\]

\[
\bar{x}_{control}=16.9\text{ cm}
\]

Hay una diferencia de:

\[
18.4-16.9=1.5\text{ cm}
\]

\pause

Pero...

\begin{alertblock}{Pregunta estadística}
¿Una diferencia de 1.5 cm es suficientemente grande como para ser difícil
de explicar por variación aleatoria?
\end{alertblock}

\end{frame}

% ------------------------------------------------------------
\begin{frame}{Variabilidad: una característica de los datos biológicos}

La biología presenta variabilidad debido a:

\begin{itemize}
    \item diferencias genéticas,
    \item edad y estado fisiológico,
    \item ambiente,
    \item error de medición,
    \item heterogeneidad espacial,
    \item variación temporal,
    \item factores no controlados.
\end{itemize}

\vspace{0.3cm}

Por eso:

\[
\text{diferencia observada}
=
\text{señal}
+
\text{variación}
\]

\pause

La estadística intenta cuantificar cuánto de la diferencia observada
es compatible con la variabilidad esperada.

\end{frame}

% ------------------------------------------------------------
\begin{frame}{La lógica de una prueba de hipótesis}

\begin{enumerate}
    \item Planteamos una hipótesis sobre la población.
    \item Recolectamos una muestra.
    \item Calculamos un estadístico.
    \item Determinamos qué tan compatible es el estadístico con $H_0$.
    \item Tomamos una decisión basada en una regla previamente definida.
\end{enumerate}

\vspace{0.4cm}

\begin{center}
\Large
\[
\boxed{
H_0
\longrightarrow
\text{datos}
\longrightarrow
\text{estadístico}
\longrightarrow
p
}
\]
\end{center}

\end{frame}

% ------------------------------------------------------------
\begin{frame}{Hipótesis nula y alternativa}

\begin{block}{Hipótesis nula ($H_0$)}
Representa generalmente una ausencia de diferencia, asociación o efecto
bajo el modelo planteado.
\end{block}

\begin{block}{Hipótesis alternativa ($H_1$ o $H_A$)}
Representa el patrón que queremos evaluar como evidencia frente a $H_0$.
\end{block}

\vspace{0.3cm}

Ejemplo:

\[
H_0:\mu_{fert}=\mu_{control}
\]

\[
H_A:\mu_{fert}\neq\mu_{control}
\]

\end{frame}

% ------------------------------------------------------------
\begin{frame}{Hipótesis unilaterales y bilaterales}

### Bilateral

\[
H_A:\mu\neq\mu_0
\]

Pregunta:

\begin{quote}
¿La media es diferente?
\end{quote}

### Unilateral

\[
H_A:\mu>\mu_0
\]

Pregunta:

\begin{quote}
¿La media es mayor?
\end{quote}

También:

\[
H_A:\mu<\mu_0
\]

\begin{alertblock}{Importante}
La dirección de la hipótesis debe estar justificada por la pregunta
científica
y definida antes de observar los resultados.
\end{alertblock}

\end{frame}

% ------------------------------------------------------------
\begin{frame}{Ejemplo: efecto de un contaminante}

Queremos estudiar si una exposición a un contaminante afecta
la tasa fotosintética.

Variable:

\[
Y = \text{tasa fotosintética}
\]

Pregunta:

\begin{quote}
¿Las plantas expuestas presentan una tasa fotosintética diferente
a las plantas control?
\end{quote}

Entonces:

\[
H_0:\mu_E=\mu_C
\]

\[
H_A:\mu_E\neq\mu_C
\]

\end{frame}

% ------------------------------------------------------------
\begin{frame}{De la pregunta biológica a la pregunta estadística}

\begin{center}
\begin{tabular}{p{0.42\textwidth} p{0.45\textwidth}}
\toprule
\textbf{Pregunta biológica} & \textbf{Hipótesis estadística}\\
\midrule
¿El tratamiento cambia el peso? &
$\mu_T \neq \mu_C$\\[0.2cm]

¿El fertilizante aumenta el crecimiento? &
$\mu_F > \mu_C$\\[0.2cm]

¿Dos hábitats tienen distinta abundancia? &
$\mu_1 \neq \mu_2$\\[0.2cm]

¿La temperatura está asociada con la respiración? &
$\rho \neq 0$\\
\bottomrule
\end{tabular}
\end{center}

\end{frame}

% ------------------------------------------------------------
\begin{frame}{El estadístico de prueba}

Una prueba transforma los datos en un número que mide la discrepancia
entre lo observado y lo esperado bajo $H_0$.

Ejemplo: prueba $t$

\[
t=
\frac{\bar{x}-\mu_0}
{s/\sqrt{n}}
\]

donde:

\begin{itemize}
    \item $\bar{x}$ = media observada,
    \item $\mu_0$ = media bajo $H_0$,
    \item $s$ = desviación estándar,
    \item $n$ = tamaño de muestra.
\end{itemize}

\end{frame}

% ------------------------------------------------------------
\begin{frame}{¿Qué significa el estadístico?}

Supongamos:

\[
\bar{x}=12.4
\]

\[
\mu_0=10
\]

\[
s=4
\]

\[
n=16
\]

Entonces:

\[
t=
\frac{12.4-10}{4/\sqrt{16}}
=
\frac{2.4}{1}
=
2.4
\]

La diferencia observada equivale a aproximadamente
$2.4$ errores estándar respecto al valor planteado por $H_0$.

\end{frame}

% ------------------------------------------------------------
\begin{frame}{¿Qué es el valor-$p$?}

El valor-$p$ responde a una pregunta específica:

\begin{block}
Si $H_0$ fuera cierta y el modelo estadístico fuera apropiado,
¿qué tan compatible sería observar un resultado como el obtenido
(o más extremo)?
\end{block}

\vspace{0.4cm}

Un valor-$p$ pequeño indica que los datos son poco compatibles con
$H_0$ bajo el modelo utilizado.

\end{frame}

% ------------------------------------------------------------
\begin{frame}{Lo que un valor-$p$ NO significa}

Si:

\[
p=0.03
\]

\textbf{NO significa:}

\begin{itemize}
    \item ``Hay 3\% de probabilidad de que $H_0$ sea cierta''.
    \item ``Hay 97\% de probabilidad de que $H_A$ sea cierta''.
    \item ``El efecto es grande''.
    \item ``El resultado es biológicamente importante''.
    \item ``El experimento está probado''.
\end{itemize}

\pause

\begin{alertblock}{Recordatorio}
El valor-$p$ se calcula suponiendo un modelo bajo $H_0$.
No es una probabilidad de que una hipótesis sea verdadera.
\end{alertblock}

\end{frame}

% ------------------------------------------------------------
\begin{frame}{Nivel de significancia $\alpha$}

Antes de analizar los resultados podemos definir:

\[
\alpha=0.05
\]

Regla convencional:

\[
p<\alpha
\quad\Rightarrow\quad
\text{rechazar }H_0
\]

\[
p\geq\alpha
\quad\Rightarrow\quad
\text{no rechazar }H_0
\]

\begin{alertblock}{Lenguaje importante}
``No rechazar $H_0$'' no significa ``demostrar que $H_0$ es verdadera''.
\end{alertblock}

\end{frame}

% ------------------------------------------------------------
\begin{frame}{Errores tipo I y tipo II}

Tenemos dos posibles errores:

\begin{center}
\begin{tabular}{c|cc}
& $H_0$ verdadera & $H_0$ falsa\\
\hline
Rechazar $H_0$ &
\textbf{Error tipo I} &
Decisión correcta\\
No rechazar $H_0$ &
Decisión correcta &
\textbf{Error tipo II}
\end{tabular}
\end{center}

\vspace{0.3cm}

\[
P(\text{Error tipo I})=\alpha
\]

\[
P(\text{Error tipo II})=\beta
\]

\end{frame}

% ------------------------------------------------------------
\begin{frame}{Ejemplo biológico: falso positivo}

Supongamos que evaluamos un nuevo tratamiento contra una enfermedad vegetal.

\[
H_0:\text{el tratamiento no tiene efecto}
\]

Un error tipo I ocurre si:

\[
\text{rechazamos }H_0
\]

cuando realmente el tratamiento no tiene efecto.

\begin{block}{Interpretación}
Concluir que existe un efecto cuando, bajo el modelo, realmente
no existe.
\end{block}

\end{frame}

% ------------------------------------------------------------
\begin{frame}{Ejemplo biológico: falso negativo}

Ahora supongamos que el tratamiento realmente funciona.

Pero nuestro estudio no encuentra evidencia suficiente para rechazar $H_0$.

Esto corresponde a un:

\[
\boxed{\text{Error tipo II}}
\]

Puede ocurrir cuando:

\begin{itemize}
    \item $n$ es pequeño,
    \item existe mucha variabilidad,
    \item el efecto es pequeño,
    \item las mediciones son poco precisas.
\end{itemize}

\end{frame}

% ------------------------------------------------------------
\begin{frame}{Potencia estadística}

La potencia es:

\[
\boxed{\text{Potencia}=1-\beta}
\]

Representa la probabilidad de detectar un efecto de una magnitud
determinada cuando ese efecto realmente existe, bajo las condiciones
del diseño y del análisis.

La potencia aumenta generalmente con:

\begin{itemize}
    \item mayor tamaño de muestra,
    \item mayor tamaño del efecto,
    \item menor variabilidad,
    \item mediciones más precisas.
\end{itemize}

\end{frame}

% ------------------------------------------------------------
\begin{frame}{Ejemplo conceptual de potencia}

Supongamos que queremos detectar una diferencia de:

\[
\Delta=2\text{ cm}
\]

entre dos tratamientos.

\begin{itemize}
    \item Con $n=5$ por grupo puede ser difícil.
    \item Con $n=20$ por grupo tenemos más información.
    \item Con $n=100$ por grupo podemos detectar diferencias pequeñas.
\end{itemize}

\begin{alertblock}{Idea clave}
Un resultado no significativo puede deberse a ausencia de efecto,
pero también a un estudio con poca capacidad para detectarlo.
\end{alertblock}

\end{frame}

% ------------------------------------------------------------
\begin{frame}{Tamaño del efecto}

La significancia estadística no mide directamente la magnitud del efecto.

Ejemplo:

\[
\Delta = 0.2\text{ cm}
\]

puede ser estadísticamente significativo con una muestra enorme.

Mientras que:

\[
\Delta = 5\text{ cm}
\]

puede no ser significativo con una muestra muy pequeña.

\begin{block}{En biología}
Debemos reportar tanto la evidencia estadística como la magnitud
y la incertidumbre del efecto.
\end{block}

\end{frame}

% ------------------------------------------------------------
\begin{frame}{Intervalos de confianza}

Un intervalo de confianza proporciona un rango de valores
compatibles con los datos y el procedimiento estadístico.

Ejemplo:

\[
\widehat{\Delta}=2.4\text{ cm}
\]

\[
IC_{95\%}=[0.8,\;4.0]\text{ cm}
\]

Podemos comunicar:

\begin{quote}
La diferencia estimada fue de 2.4 cm, con un intervalo de confianza
del 95\% de 0.8 a 4.0 cm.
\end{quote}

\end{frame}

% ------------------------------------------------------------
\begin{frame}{Pruebas de hipótesis comunes}

\begin{center}
\begin{tabular}{lll}
\toprule
\textbf{Pregunta} & \textbf{Prueba frecuente} & \textbf{Alternativa}\\
\midrule
Una media vs. valor & $t$ de una muestra & Wilcoxon\\
Dos grupos & $t$ de dos muestras & Mann--Whitney\\
Antes/después & $t$ pareada & Wilcoxon pareada\\
Más de dos grupos & ANOVA & Kruskal--Wallis\\
Dos variables categóricas & $\chi^2$ & Fisher\\
Asociación lineal & correlación $t$ & Spearman\\
\bottomrule
\end{tabular}
\end{center}

\end{frame}

% ------------------------------------------------------------
\begin{frame}{Ejemplo 1: una muestra}

Pregunta:

\begin{quote}
¿La longitud promedio de las hojas de una población es diferente
de 10 cm?
\end{quote}

Hipótesis:

\[
H_0:\mu=10
\]

\[
H_A:\mu\neq10
\]

Si las condiciones son razonables:

\[
\text{prueba }t\text{ de una muestra}
\]

\end{frame}

% ------------------------------------------------------------
\begin{frame}{Ejemplo 2: dos grupos independientes}

Tenemos dos grupos:

\begin{itemize}
    \item Plantas control.
    \item Plantas tratadas.
\end{itemize}

Variable respuesta:

\[
Y=\text{biomasa seca}
\]

Pregunta:

\[
\mu_T-\mu_C\neq0?
\]

Una opción frecuente:

\[
\boxed{\text{prueba }t\text{ para dos muestras}}
\]

\end{frame}

% ------------------------------------------------------------
\begin{frame}{Prueba $t$: idea intuitiva}

La prueba compara:

\[
\text{diferencia entre medias}
\]

con:

\[
\text{variabilidad esperada de esa diferencia}
\]

En términos intuitivos:

\[
t=
\frac{\text{señal}}
{\text{ruido}}
\]

\pause

Cuanto mayor sea la diferencia respecto a su incertidumbre,
más evidencia puede existir contra $H_0$.

\end{frame}

% ------------------------------------------------------------
\begin{frame}{Ejemplo 3: datos pareados}

Medimos la presión arterial de animales:

\[
\text{antes}
\qquad\text{y}\qquad
\text{después}
\]

Los datos están relacionados porque:

\[
\boxed{\text{cada animal aporta dos mediciones}}
\]

Podemos trabajar con las diferencias:

\[
D_i=Después_i-Antes_i
\]

y probar:

\[
H_0:\mu_D=0
\]

\end{frame}

% ------------------------------------------------------------
\begin{frame}{Ejemplo 4: más de dos grupos}

Queremos comparar la biomasa de plantas bajo:

\begin{itemize}
    \item Control
    \item Fertilizante A
    \item Fertilizante B
    \item Fertilizante C
\end{itemize}

Una prueba ANOVA puede evaluar:

\[
H_0:
\mu_C=\mu_A=\mu_B=\mu_C
\]

frente a:

\[
H_A:\text{al menos una media difiere}
\]

\end{frame}

% ------------------------------------------------------------
\begin{frame}{ANOVA: una advertencia importante}

Un resultado significativo de ANOVA nos dice:

\[
\boxed{\text{no todas las medias son iguales}}
\]

Pero no identifica automáticamente:

\[
\text{qué grupos difieren entre sí}
\]

Para eso se requieren comparaciones posteriores apropiadas
(post-hoc), teniendo en cuenta la multiplicidad.

\end{frame}

% ------------------------------------------------------------
\begin{frame}{Pruebas no paramétricas}

Cuando las condiciones de una prueba paramétrica no son razonables,
pueden existir alternativas.

Ejemplos:

\begin{center}
\begin{tabular}{ll}
\toprule
Paramétrica & Alternativa frecuente\\
\midrule
$t$ independiente & Mann--Whitney\\
$t$ pareada & Wilcoxon\\
ANOVA & Kruskal--Wallis\\
Correlación de Pearson & Spearman\\
\bottomrule
\end{tabular}
\end{center}

\begin{alertblock}{Pero}
``No paramétrica'' no significa ``sin supuestos''.
\end{alertblock}

\end{frame}

% ------------------------------------------------------------
\begin{frame}{Supuestos: ¿por qué importan?}

Una prueba estadística depende de un modelo.

Por ejemplo, para una prueba $t$ pueden ser importantes:

\begin{itemize}
    \item independencia de las observaciones,
    \item estructura adecuada del diseño,
    \item distribución aproximadamente apropiada de los residuos,
    \item ausencia de valores extremos que dominen el resultado.
\end{itemize}

\begin{block}{Principio}
La prueba debe corresponder al diseño y a la estructura de los datos.
\end{block}

\end{frame}

% ------------------------------------------------------------
\begin{frame}{Independencia: un problema muy común}

Supongamos que medimos:

\[
10\text{ hojas por planta}
\]

de:

\[
5\text{ plantas}
\]

Tenemos 50 mediciones, pero no necesariamente:

\[
n=50
\]

observaciones independientes.

Las hojas de una misma planta pueden estar correlacionadas.

\begin{alertblock}{Pseudorreplicación}
Tratar mediciones no independientes como réplicas independientes
puede producir inferencias engañosas.
\end{alertblock}

\end{frame}

% ------------------------------------------------------------
\begin{frame}{¿Qué hacer antes de una prueba?}

Antes de calcular un valor-$p$:

\begin{enumerate}
    \item Definir la pregunta científica.
    \item Identificar la unidad experimental.
    \item Revisar el diseño.
    \item Explorar los datos.
    \item Identificar variables y grupos.
    \item Considerar supuestos.
    \item Elegir el modelo/prueba.
\end{enumerate}

\end{frame}

% ------------------------------------------------------------
\begin{frame}{Flujo de trabajo}

\begin{center}

\begin{tikzpicture}[node distance=1.3cm]

\node[draw, rounded corners, fill=green!10] (q) {Pregunta biológica};
\node[draw, rounded corners, fill=blue!10, right=of q] (d) {Diseño};
\node[draw, rounded corners, fill=yellow!15, right=of d] (e) {Exploración};
\node[draw, rounded corners, fill=orange!15, right=of e] (t) {Prueba/modelo};
\node[draw, rounded corners, fill=red!10, right=of t] (i) {Interpretación};

\draw[->, thick] (q) -- (d);
\draw[->, thick] (d) -- (e);
\draw[->, thick] (e) -- (t);
\draw[->, thick] (t) -- (i);

\end{tikzpicture}

\end{center}

\vspace{0.5cm}

\centering
\textbf{La prueba estadística no sustituye el razonamiento científico.}

\end{frame}

% ------------------------------------------------------------
\begin{frame}[fragile]{Ejemplo en R}

Supongamos un conjunto de datos:

\begin{lstlisting}[language=R]
control <- c(12.1, 11.4, 13.2, 12.8, 11.9)
tratado <- c(15.0, 14.2, 16.1, 15.4, 14.8)

t.test(tratado, control)
\end{lstlisting}

La función produce, entre otros:

\begin{itemize}
    \item diferencia estimada,
    \item estadístico $t$,
    \item grados de libertad,
    \item valor-$p$,
    \item intervalo de confianza.
\end{itemize}

\end{frame}

% ------------------------------------------------------------
\begin{frame}[fragile]{Ejemplo en Python}

Con \texttt{scipy}:

\begin{lstlisting}[language=Python]
from scipy.stats import ttest_ind

control = [12.1, 11.4, 13.2, 12.8, 11.9]
tratado = [15.0, 14.2, 16.1, 15.4, 14.8]

resultado = ttest_ind(tratado, control)

print(resultado)
\end{lstlisting}

\end{frame}

% ------------------------------------------------------------
\begin{frame}{Interpretación de un resultado}

Supongamos que obtenemos:

\[
t=3.12
\]

\[
p=0.014
\]

\[
IC_{95\%}=[0.55,\;3.65]\text{ g}
\]

Una forma adecuada de comunicarlo:

\begin{block}
La biomasa media fue mayor en el grupo tratado que en el control,
con una diferencia estimada de 2.1 g (IC 95\%: 0.55--3.65 g;
$p=0.014$).
\end{block}

\end{frame}

% ------------------------------------------------------------
\begin{frame}{¿Significativo = importante?}

No necesariamente.

Consideremos:

\[
p<0.05
\]

Esto proporciona evidencia estadística contra $H_0$ según el análisis.

Pero todavía debemos preguntar:

\begin{itemize}
    \item ¿Qué tan grande es el efecto?
    \item ¿Es biológicamente relevante?
    \item ¿Cuál es su incertidumbre?
    \item ¿El diseño permite responder la pregunta?
    \item ¿Existen posibles sesgos?
\end{itemize}

\end{frame}

% ------------------------------------------------------------
\begin{frame}{Problema de comparaciones múltiples}

Supongamos que hacemos 20 pruebas independientes con:

\[
\alpha=0.05
\]

Incluso si todas las hipótesis nulas fueran ciertas,
la probabilidad de obtener al menos un resultado significativo puede ser
alta.

\[
P(\text{al menos un falso positivo})
=
1-(1-0.05)^{20}
\]

\[
\approx 0.64
\]

\begin{alertblock}{Idea}
Muchas pruebas requieren considerar el problema de multiplicidad.
\end{alertblock}

\end{frame}

% ------------------------------------------------------------
\begin{frame}{Correcciones por comparaciones múltiples}

Algunas estrategias:

\begin{itemize}
    \item Bonferroni
    \item Holm
    \item Benjamini--Hochberg
    \item Modelos que incorporan simultáneamente varios factores
\end{itemize}

La elección depende de:

\begin{itemize}
    \item objetivo inferencial,
    \item número de comparaciones,
    \item estructura del estudio,
    \item control deseado del error.
\end{itemize}

\end{frame}

% ------------------------------------------------------------
\begin{frame}{Ejercicio de discusión}

Un investigador compara la abundancia de aves entre dos bosques.

Obtiene:

\[
p=0.08
\]

Pregunta:

\begin{enumerate}
    \item ¿Podemos decir que ``no existe diferencia''?
    \item ¿Qué información adicional necesitamos?
    \item ¿Qué papel juega el tamaño de muestra?
    \item ¿Qué papel juega el tamaño del efecto?
    \item ¿Qué nos diría un intervalo de confianza?
\end{enumerate}

\vspace{0.3cm}

\pause

\begin{block}{Discusión}
No rechazar $H_0$ no equivale a demostrar que las poblaciones son idénticas.
\end{block}

\end{frame}

% ------------------------------------------------------------
\begin{frame}{Ejercicio: elegir la prueba}

Seleccione una estrategia para cada situación:

\begin{enumerate}
    \item Comparar crecimiento entre tratamiento y control.
    \item Comparar crecimiento antes y después en los mismos individuos.
    \item Comparar biomasa entre cuatro tratamientos.
    \item Evaluar asociación entre temperatura y tasa metabólica.
    \item Evaluar asociación entre dos variables categóricas.
\end{enumerate}

\vspace{0.3cm}

\pause

\textbf{Discusión:}

\begin{itemize}
    \item $t$ independiente
    \item $t$ pareada
    \item ANOVA
    \item correlación/regresión
    \item $\chi^2$ o Fisher
\end{itemize}

\end{frame}

% ------------------------------------------------------------
\begin{frame}{Ejercicio integrador}

Tenemos datos de crecimiento de plantas.

\begin{center}
\begin{tabular}{ccc}
\toprule
Planta & Control & Fertilizante\\
\midrule
1 & 10.2 & 12.1\\
2 & 11.4 & 12.8\\
3 & 9.8 & 11.9\\
4 & 10.7 & 13.0\\
5 & 10.9 & 12.5\\
6 & 11.2 & 13.1\\
\bottomrule
\end{tabular}
\end{center}

Preguntas:

\begin{enumerate}
    \item ¿Cuál es la unidad experimental?
    \item ¿Cuál es $H_0$?
    \item ¿Cuál es $H_A$?
    \item ¿Qué prueba usarían?
    \item ¿Qué supuestos revisarían?
    \item ¿Qué resultado reportarían?
\end{enumerate}

\end{frame}

% ------------------------------------------------------------
\begin{frame}{Ejercicio integrador: razonamiento}

Una posible formulación:

\[
H_0:\mu_F-\mu_C=0
\]

\[
H_A:\mu_F-\mu_C\neq0
\]

Si los grupos son independientes y los supuestos son razonables:

\[
\boxed{\text{prueba }t\text{ de dos muestras}}
\]

Pero la prueba final depende del diseño real.

\end{frame}

% ------------------------------------------------------------
\begin{frame}{Errores frecuentes}

\begin{itemize}
    \item Elegir la prueba mirando únicamente el valor-$p$.
    \item Interpretar $p$ como la probabilidad de que $H_0$ sea cierta.
    \item Confundir ``no significativo'' con ``sin efecto''.
    \item Ignorar la independencia.
    \item Hacer muchas pruebas sin considerar multiplicidad.
    \item Reportar únicamente el valor-$p$.
    \item Ignorar el tamaño del efecto.
    \item Convertir automáticamente resultados estadísticos en conclusiones
    biológicas.
\end{itemize}

\end{frame}

% ------------------------------------------------------------
\begin{frame}{¿Cómo reportar una prueba de hipótesis?}

Una buena comunicación incluye:

\begin{enumerate}
    \item Qué se comparó.
    \item Qué efecto se estimó.
    \item Tamaño del efecto.
    \item Intervalo de confianza.
    \item Prueba/modelo utilizado.
    \item Valor-$p$.
    \item Tamaño de muestra.
    \item Información relevante sobre el diseño.
\end{enumerate}

\vspace{0.3cm}

\begin{block}{Ejemplo}
El tratamiento aumentó la biomasa en 2.1 g
(IC 95\%: 0.55--3.65; $t=3.12$, $p=0.014$).
\end{block}

\end{frame}

% ------------------------------------------------------------
\begin{frame}{Mapa mental}

\begin{center}
\Large

\[
\boxed{\text{Pregunta biológica}}
\]

$\downarrow$

\[
\boxed{\text{Diseño y datos}}
\]

$\downarrow$

\[
\boxed{H_0\quad vs.\quad H_A}
\]

$\downarrow$

\[
\boxed{\text{Estadístico + incertidumbre}}
\]

$\downarrow$

\[
\boxed{\text{Valor-}p}
\]

$\downarrow$

\[
\boxed{\text{Tamaño del efecto}}
\]

$\downarrow$

\[
\boxed{\text{Interpretación biológica}}
\]

\end{center}

\end{frame}

% ------------------------------------------------------------
\begin{frame}{Resumen}

\begin{itemize}
    \item Las pruebas de hipótesis permiten cuantificar evidencia
    frente a una hipótesis nula bajo un modelo estadístico.
    \item El valor-$p$ no es la probabilidad de que una hipótesis sea
    verdadera.
    \item Debemos distinguir errores tipo I y tipo II.
    \item La potencia depende, entre otras cosas, del tamaño de muestra,
    efecto y variabilidad.
    \item La significancia estadística no equivale a importancia biológica.
    \item El diseño experimental y la independencia son fundamentales.
    \item Debemos reportar efectos e incertidumbre, no únicamente
    valores-$p$.
\end{itemize}

\end{frame}

% ------------------------------------------------------------
\begin{frame}{Preguntas para cerrar la clase}

\begin{enumerate}
    \item ¿Qué pregunta responde realmente un valor-$p$?
    \item ¿Cuál es la diferencia entre ``no rechazar $H_0$''
    y ``aceptar $H_0$''?
    \item ¿Qué es un error tipo I?
    \item ¿Qué factores afectan la potencia?
    \item ¿Por qué el tamaño del efecto es importante?
    \item ¿Por qué la independencia importa tanto en biología?
\end{enumerate}

\vspace{0.5cm}

\centering
\Large
\textbf{La estadística debe ayudar a responder preguntas biológicas,
no reemplazarlas.}

\end{frame}

% ------------------------------------------------------------
\begin{frame}{Actividad posterior}

\begin{block}{Tarea}
Tome un conjunto de datos biológicos y construya un análisis que incluya:
\end{block}

\begin{enumerate}
    \item Pregunta científica.
    \item Variable respuesta.
    \item Variable explicativa.
    \item Unidad experimental.
    \item $H_0$ y $H_A$.
    \item Visualización exploratoria.
    \item Prueba o modelo seleccionado.
    \item Tamaño del efecto.
    \item Intervalo de confianza.
    \item Valor-$p$.
    \item Interpretación biológica.
\end{enumerate}

\end{frame}

% ------------------------------------------------------------
\begin{frame}
\centering
\Huge
\textcolor{bioGreen}{Gracias}

\vspace{0.5cm}

\Large
¿Preguntas?
\end{frame}

% ============================================================
\end{document}
% ============================================================
